\section{Hyperbolic and projective views of constant torque}
\label{sec:hyperbolic}

\subsection{A product suggests a reciprocal motion}

For the elementary curve \(xy=c>0\), replacing
\(x\) by \(e^\alpha x\) and \(y\) by \(e^{-\alpha}y\) leaves the product
unchanged. This is how hyperbolic coordinates are discovered: one factor grows
exactly as the other shrinks. In the torque plane, rotate by \(45^\circ\):
\begin{equation}
 \gls{sym:uvirt}=\frac{\gls{sym:zpsi}+\gls{sym:zq}}{\sqrt2},\qquad
 \gls{sym:vvirt}=\frac{\gls{sym:zpsi}-\gls{sym:zq}}{\sqrt2}.
 \label{eq:uv-rotation}
\end{equation}
Then
\begin{equation}
 \gls{sym:torque}=\frac{\gls{sym:kappa}}2(u^2-v^2),\qquad
 \gls{sym:pcu}=u^2+v^2+z_0^2.
 \label{eq:signature-form}
\end{equation}
No invertible linear transformation can turn \(u^2-v^2\) into a single
positive square: the saddle has one positive, one negative, and one zero
direction. Sylvester's law of inertia formalizes this geometric fact---a
congruence change preserves the counts of positive, negative, and zero
directions \cite{horn2013}.

For positive torque define
\begin{equation}
 \gls{sym:it}=\sqrt{\frac{2\gls{sym:torque}}{\gls{sym:kappa}}},\qquad
 u=i_T\cosh\gls{sym:alpha},\quad v=i_T\sinh\gls{sym:alpha}.
 \label{eq:hyperbolic-coordinates}
\end{equation}
Because \(\cosh^2\alpha-\sinh^2\alpha=1\), torque depends only on the virtual
amplitude \(i_T\), whereas
\begin{equation}
 \gls{sym:pcu}=i_T^2\cosh(2\alpha)+z_0^2.
 \label{eq:hyperbolic-loss}
\end{equation}
The inequality \(\cosh(2\alpha)\geq1\) proves without differentiation that
balanced factors, \(\alpha=0\), minimize loss when voltage is inactive. The
constant-torque hyperbola is an orbit of the reciprocal scaling above, which
is mathematically the same one-parameter group as a planar hyperbolic
rotation. The Lorentz analogy adds no machine physics and is not needed beyond
this invariant-product observation.

\begin{figure}[tb]
\centering
\begin{tikzpicture}
\begin{axis}[axis main,width=.76\linewidth,height=6.1cm,
 axis equal image,xlabel={$u$},ylabel={$v$},xmin=-3,xmax=3,ymin=-2.2,ymax=2.2]
 \addplot[torque curve,domain=1:3,samples=100]{sqrt(x^2-1)};
 \addplot[torque curve,domain=1:3,samples=100]{-sqrt(x^2-1)};
 \addplot[torque curve,domain=-3:-1,samples=100]{sqrt(x^2-1)};
 \addplot[torque curve,domain=-3:-1,samples=100]{-sqrt(x^2-1)};
 \addplot[loss curve,domain=0:360,samples=121] ({1.75*cos(x)},{1.75*sin(x)});
 \addplot[only marks,mark=*,ForestGreen] coordinates {(1,0)};
 \node[annotation,anchor=west] at (axis cs:1.08,.05){$\alpha=0$};
 \node[annotation] at (axis cs:2.25,1.55){same product,\\more loss};
\end{axis}
\end{tikzpicture}
\caption{A constant-torque hyperbola is parameterized by a hyperbolic angle.
The smallest loss circle touches it at balanced factors, $\alpha=0$.}
\label{fig:hyperbolic}
\end{figure}


\subsection{Remove radius with ratios}

On a positive-torque chart where \(z_q\neq0\), define
\begin{equation}
 \gls{sym:eta}=\frac{\gls{sym:zpsi}}{\gls{sym:zq}}>0,
 \qquad \gls{sym:chi}=\frac{\gls{sym:znull}}{\gls{sym:zq}}.
 \label{eq:projective-coordinates}
\end{equation}
The torque equation fixes
\(z_q^2=M/(\kappa\eta)\). Substitution into the loss sphere gives
\begin{equation}
 \gls{sym:pcu}=\frac{M}{\kappa}\gls{sym:jproj}(\eta,\chi),
 \qquad
 \boxed{\gls{sym:jproj}(\eta,\chi)=
 \eta+\frac{1+\chi^2}{\eta}.}
 \label{eq:J-projective}
\end{equation}
Setting \(J=c\) and multiplying by \(\eta\) yields
\begin{equation}
 \boxed{\left(\eta-\frac c2\right)^2+\chi^2
 =\frac{c^2}{4}-1.}
 \label{eq:J-circles}
\end{equation}
The circles are not a plotting coincidence. Dividing a spherical loss by the
bilinear torque product produces the reciprocal term \(1/\eta\); completing
the square returns Euclidean circles in the ratio chart. Since
\(\eta+1/\eta\geq2\), the level \(c=2\) degenerates to the single point
\((\eta,\chi)=(1,0)\): balanced active flux, balanced quadrature factor, and
no torque-neutral current. Larger circles represent increasingly inefficient
allocations. For prescribed torque the loss ceiling is the outer circle
\begin{equation}
 c_{\max}=\frac{\kappa\gls{sym:pcumax}}{|M|}.
 \label{eq:outer-loss-circle}
\end{equation}

Let \(\bar{\matr Q}_U=[q_{jk}]\) be the voltage matrix in \(\vect z\)
coordinates and \(\vect d=[\eta,1,\chi]^{\mathsf T}\). The voltage condition
becomes
\begin{equation}
 \vect d^{\mathsf T}\bar{\matr Q}_U\vect d
 \leq\frac{\kappa\eta\gls{sym:umax}^2}{|M|}.
 \label{eq:projective-voltage}
\end{equation}
This is a conic in \((\eta,\chi)\). For \(R_s>0\), \(L_{\mathrm{exc}}>0\),
and \(\omega_e\neq0\), its quadratic block in \((\eta,\chi)\) is positive
definite and the admissible set is ellipse-like when nonempty. At
\(R_s=0\), the two active-flux-plane voltage columns become collinear and the
conic can degenerate. Minimum loss is now literally the smallest \(J\)-circle
that meets the voltage-admissible conic. Maximum torque uses the same contact
picture while the outer loss circle is fixed and the torque parameter changes.

\begin{figure}[tb]
\centering
\begin{tikzpicture}
\begin{axis}[axis main,width=.78\linewidth,height=6.3cm,axis equal image,
 xlabel={$\eta$},ylabel={$\chi$},xmin=0,xmax=5.1,ymin=-2.1,ymax=2.1]
 \addplot[loss curve,domain=0:360,samples=121]
 ({1+0*cos(x)},{0+0*sin(x)});
 \addplot[loss curve,MidnightBlue!65,domain=0:360,samples=121]
 ({1.25+.75*cos(x)},{.75*sin(x)});
 \addplot[loss curve,MidnightBlue!40,domain=0:360,samples=121]
 ({1.75+1.436*cos(x)},{1.436*sin(x)});
 \addplot[voltage curve,domain=0:360,samples=121]
 ({3.15+1.05*cos(x)*cos(18)-.62*sin(x)*sin(18)},
  {.15+1.05*cos(x)*sin(18)+.62*sin(x)*cos(18)});
 \addplot[only marks,mark=*,ForestGreen] coordinates {(1,0)};
 \node[annotation,anchor=west] at (axis cs:1.05,.05){$J=2$};
 \node[annotation] at (axis cs:3.55,-1.25){voltage conic};
\end{axis}
\end{tikzpicture}
\caption{Projective loss levels are circles with moving centers. The
unconstrained optimum is the degenerate point $(1,0)$; voltage forces the first
contact with its transformed conic.}
\label{fig:projective}
\end{figure}


Finally, the hyperbolic and projective parameters are the same allocation
written linearly and logarithmically. From \cref{eq:uv-rotation} and
\cref{eq:hyperbolic-coordinates},
\begin{equation}
 \eta=\frac{u+v}{u-v}=e^{2\alpha},\qquad
 \alpha=\frac12\log\eta.
 \label{eq:eta-alpha}
\end{equation}
Ratios are convenient for algebra and conics; the hyperbolic angle turns
multiplicative imbalance into an additive coordinate.
